\documentclass[10pt,a4paper]{amsart}
\usepackage[margin=19mm]{geometry}
\usepackage{amsmath,amssymb,mathtools}
\usepackage[scr=rsfs,cal=euler]{mathalfa}
\usepackage{xfrac}
\usepackage[unicode,hidelinks,bookmarks=false]{hyperref}
\newcommand{\ie}{i.e.\ }
\newcommand{\cf}{cf.\ }
\newcommand{\resp}{respectively\ }
\newcommand{\Subsection}{\subsection}
\providecommand{\nobreakdash}{\nobreak\hbox{-}}
\setlength{\emergencystretch}{2em}
\multlinegap0pt
\usepackage{mathtools}
\usepackage{amssymb}
\usepackage{xcolor}
\usepackage{float}
\usepackage{tikz}
\newtheorem{theorem}{Theorem}
\newtheorem{lemma}{Lemma}
\newcommand{\R}{\mathbb R}%
\title{Ces\`aro summability of H\"older functions and Talbot effect on rank one Riemannian symmetric spaces of compact type}
\author{Utsav Dewan}
\date{Source: arXiv:2604.02305v1 (CC BY 4.0)}
\begin{document}
\maketitle

\noindent\emph{Source and licence.} Ces\`aro summability of H\"older functions and Talbot effect on rank one Riemannian symmetric spaces of compact type. Version arXiv:2604.02305v1 (CC BY 4.0), distributed by arXiv under the Creative Commons Attribution 4.0 International licence, http://creativecommons.org/licenses/by/4.0/, as recorded in the arXiv licence metadata for this exact version and shown on the abstract page https://arxiv.org/abs/2604.02305v1. Full licence notice: \texttt{licenses/arxiv-2604.02305-CC-BY-4.0.txt}.\par\smallskip

\noindent\textit{Editorial scope.} Counted unit: the article's theorem thm\_gmt (source0.tex lines 739--746). The article's complete original proof of it is delivered with it (source0.tex lines 825--857, one \texttt{proof} environment of 33 lines, which the article places away from the statement). No mathematics is written, repaired or re-derived: the statement and the proof are the article's own lines, in the article's own notation, and no retained line is substituted or re-flowed.  Retained uncounted as the source-local context the proof needs: lines 749--752; lines 763--765; lines 789--794.  Nothing was omitted from any retained span, so the package carries no omission record.  No retained line contains a citation, so the wrapper carries no bibliography and no bibliography entry.  No retained line contains a figure environment, an image-inclusion command or drawing code, so no image is delivered and the package declares no asset dependency.  Editorial formatting only, in addition: an AMS \texttt{amsart} preamble that restates the article's own declarations and macros used by the retained lines, namely the environments \texttt{theorem}, \texttt{lemma} on separate counters, together with the standard packages those lines need.  The retained environments print in this document as Theorem 1, Lemma 1 in order of appearance, whereas the article prints the same items with the numbering of its own full document.  The complete native source of the article as distributed in the arXiv e-print for this exact version is bundled unchanged as \texttt{sources/197726/source0.tex}.
\begin{lemma} \label{gmt_lemma3}
Let $E$ be a bounded subset of $\R^m$ and $f \in C^\gamma(E),\: \gamma \in (0,1)$, be real-valued. Then 
$$\overline{\dim}_M\left(Graph(f)\right) \le (m+1)-\gamma\:.$$
\end{lemma}

\begin{lemma} \label{gmt_lemma1}
Let $\varphi$ be a bijective bi-Lipschitz map from $(\mathcal{X}_1,d_1) \to (\mathcal{X}_2,d_2)$ where both are totally bounded metric spaces. Then $\overline{\dim}_M\left(\mathcal{X}_1\right)=\overline{\dim}_M\left(\mathcal{X}_2\right)$\:.
\end{lemma}

\begin{lemma} \label{gmt_lemma2}
Let  $\left(\mathcal{X},d\right)$ be a totally bounded metric space and $\Omega_j \subset \mathcal{X},\:1\le j \le n$. Then 
\begin{equation*}
\overline{\dim}_M\left(\bigcup_{j=1}^n\Omega_j\right) = \max_{1 \le j \le n}\overline{\dim}_M\left(\Omega_j\right)\:.
\end{equation*} 
\end{lemma}

\begin{theorem} \label{thm_gmt}
Let $(\mathcal{X},d)$ be a compact metric space such that there are finitely many subsets $\Omega_j \subset \mathcal{X},\:1 \le j \le n$ with
\begin{itemize}
\item $\mathcal{X} = \displaystyle\bigcup_{j=1}^n \Omega_j$ and
\item there are bi-Lipschitz maps $\varphi_j : \Omega_j \to \varphi_j\left(\Omega_j\right) \subset \mathbb{R}^m\:,\: 1\le j \le n$. 
\end{itemize}
Then for a real-valued $f \in C^\gamma(\mathcal{X}),\:\gamma \in (0,1)$, we have $\overline{\dim}_M\left(Graph(f)\right) \le (m+1)-\gamma$\:.
\end{theorem}

\begin{proof}[Proof of Theorem \ref{thm_gmt}]
We set for $1 \le j \le n$,
\begin{equation*}
G_j:= \left\{\left(x,f(x)\right) : x \in \Omega_j\right\} \subset Graph(f)\:.
\end{equation*}
Then
\begin{equation} \label{thm_gmt_pf_eq1}
Graph(f) = \bigcup_{j=1}^n G_j\:.
\end{equation}
Now as $\mathcal{X}$ is compact and $f$ is continuous, we have that $Graph(f)$ is a compact subset of $\mathcal{X} \times \R$ and consequently all the above subsets $G_j$ are totally bounded. 

Next for each $1 \le j \le n$, we set $\tilde{G}_j$ to be the graph of  $f \circ \varphi^{-1}_j$ on $\varphi_j\left(\Omega_j\right)$, that is,
\begin{equation*}
\tilde{G}_j:=\left\{\left(y,f \circ \varphi^{-1}_j(y)\right) : y \in \varphi_j\left(\Omega_j\right)\right\}\:.
\end{equation*}
We next define,
\begin{eqnarray*}
\Phi_j : G_j &\to & \tilde{G}_j \\
\left(x,f(x)\right) &\mapsto & \left(\varphi_j(x),f(x)\right)\:.
\end{eqnarray*}
As $\varphi_j$ is a bi-Lipschitz map from $\left(\Omega_j,d\right)$ onto $\left(\varphi_j\left(\Omega_j\right),d_{\R^m}\right)$ ($d_{\R^m}$ denotes the canonical metric on $\R^m$), it follows that $\Phi_j$ is a bi-Lipschitz map from $\left(G_j,\:d_{\mathcal{X} \times \R}\right)$ onto  $\left(\tilde{G}_j,\:d_{\R^{m+1}}\right)$, where $d_{\mathcal{X} \times \R}$ denotes the restriction of the product metric on $\mathcal{X} \times \R$ to $G_j$. Thus $\tilde{G}_j$ being the image of a totally bounded set $G_j$ under $\Phi_j$ is also totally bounded and moreover, by Lemma \ref{gmt_lemma1},
\begin{equation} \label{thm_gmt_pf_eq2}
\overline{\dim}_M\left(G_j\right)=\overline{\dim}_M\left(\tilde{G}_j\right)\:.
\end{equation}
Next we note that as $f \in C^\gamma(\mathcal{X})$, we also have  $f \circ \varphi^{-1}_j \in C^\gamma(\varphi_j\left(\Omega_j\right))$ and hence, $\tilde{G}_j$ is the graph of a real-valued $C^\gamma$ function on a bounded set $\varphi_j\left(\Omega_j\right) \subset \R^m$. Then by Lemma \ref{gmt_lemma3}, we have that
\begin{equation} \label{thm_gmt_pf_eq3}
\overline{\dim}_M\left(\tilde{G}_j\right) \le (m+1)-\gamma\:.
\end{equation}
Finally, applying Lemma \ref{gmt_lemma2} on (\ref{thm_gmt_pf_eq1}) and then combining it with (\ref{thm_gmt_pf_eq2}) and (\ref{thm_gmt_pf_eq3}), we get the desired bound,
\begin{equation*}
\overline{\dim}_M\left(Graph(f)\right) =\max_{1 \le j \le n}\overline{\dim}_M\left(G_j\right)=\max_{1 \le j \le n}\overline{\dim}_M\left(\tilde{G}_j\right) \le (m+1)-\gamma\:.
\end{equation*}
\end{proof}

\end{document}
